\phantomsection\addcontentsline{toc}{section}{Internet of Things}
\ECchapter{17}{Internet of Things}{How can connected objects remain useful, secure and energy-efficient?}{ECGreen}

\ECObjectives{
\begin{itemize}
\item Describe the architecture of an IoT system from sensing to user service.
\item Compare communication technologies and explain the importance of energy management.
\item Identify reliability, interoperability, privacy and cybersecurity requirements.
\end{itemize}}

\begin{multicols}{2}
\begin{engineeringnews}
Connected sensors now monitor factories, buildings, farms and transport networks. The engineering
challenge is no longer simply to connect a device, but to keep thousands of low-power objects useful,
secure and maintainable throughout their entire service life.
\end{engineeringnews}

\begin{ecarticle}
The Internet of Things (IoT) links physical objects to digital services. A typical object contains a
sensor, a microcontroller, a communication interface and a power supply. It measures a physical
quantity, processes the signal locally and sends useful information to a gateway or cloud platform.
The platform stores data, detects abnormal behaviour and may send a command back to an actuator.

A connected temperature sensor illustrates the complete chain. The sensing element converts
temperature into an electrical signal. The microcontroller filters the measurement, adds a timestamp
and decides whether transmission is necessary. Sending every value provides detailed data but wastes
energy and network capacity. Edge processing can instead transmit only a periodic summary or an alarm
when a threshold is exceeded.

Communication technology depends on range, data rate, power and infrastructure. Bluetooth Low Energy
is suited to short-range devices. Wi-Fi offers higher throughput but usually consumes more power.
LoRaWAN can transmit small packets over several kilometres, while cellular IoT provides wide-area
coverage through an operator network. No protocol is best for every application.

Energy autonomy is often the limiting requirement. Engineers reduce consumption by using sleep modes,
low sampling rates and short radio transmissions. A battery-powered sensor may need to operate for
several years without maintenance. Energy harvesting from light, vibration or heat can extend battery life,
but the available power is small and variable.

Security must be designed from the start. Each object needs a unique identity, authenticated software
updates and encrypted communication. A compromised device can reveal private data or become an entry
point into a larger system. Finally, interoperability matters: devices from different manufacturers
must use compatible data formats so that the system can evolve without replacing every component.
\end{ecarticle}

\begin{tcolorbox}[title=\sffamily\bfseries Did You Know?,colback=yellow!10,colframe=ECOrange]
For a battery-powered object, radio transmission may consume far more energy than the measurement
itself. Sending less data is therefore often more effective than improving the sensor electronics.
\end{tcolorbox}

\begin{engineeringfocus}
\textbf{Architecture of a connected monitoring system}
\begin{center}
\resizebox{0.98\linewidth}{!}{%
\begin{tikzpicture}[font=\scriptsize,>=Latex,node distance=4mm]
\node[draw=ECGreen,fill=ECGreen!10,rounded corners,align=center] (s) {sensor and\\microcontroller};
\node[draw=ECBlue,fill=ECLightBlue,rounded corners,align=center,right=of s] (r) {low-power\\radio};
\node[draw=ECPurple,fill=ECPurple!10,rounded corners,align=center,right=of r] (g) {gateway or\\base station};
\node[draw=ECOrange,fill=ECOrange!10,rounded corners,align=center,below=of g] (p) {data platform\\and analytics};
\node[draw=ECRed,fill=ECRed!8,rounded corners,align=center,below=of r] (u) {dashboard or\\maintenance alarm};
\draw[->,thick] (s)--(r); \draw[->,thick] (r)--(g); \draw[->,thick] (g)--(p);
\draw[->,thick] (p)--(u); \draw[->,thick,dashed] (u)--(s);
\end{tikzpicture}%
}
\end{center}
Local processing reduces traffic and allows basic operation even if the cloud connection is lost.
\end{engineeringfocus}

\begin{keyfigures}
\begin{tabularx}{\linewidth}{@{}Xr@{}}
Typical sensor data rate & low\\
Critical resource & battery energy\\
Long-range low-power example & LoRaWAN\\
Short-range example & BLE\\
Essential maintenance feature & remote update
\end{tabularx}
\end{keyfigures}

\begin{tcolorbox}[title=\sffamily\bfseries Relative communication current,colback=white,colframe=ECGreen]
\begin{center}
\begin{tikzpicture}
\begin{axis}[xbar,width=0.86\linewidth,height=50mm,xlabel={Relative current},xmin=0,xmax=110,
symbolic y coords={Deep sleep,BLE advertising,LoRaWAN,Wi-Fi duty cycle,Continuous Wi-Fi},
ytick=data,grid=major,ticklabel style={font=\scriptsize},nodes near coords]
\addplot table[x=current,y=mode,col sep=comma]{data/EC17/iot_power.csv};
\end{axis}
\end{tikzpicture}
\end{center}
\footnotesize Illustrative values: actual consumption depends strongly on hardware, coverage and transmission frequency.
\end{tcolorbox}

\begin{vocabulary}
\begin{tabularx}{\linewidth}{@{}lX@{}}
connected object & objet connecté\\
gateway & passerelle\\
throughput & débit\\
sleep mode & mode veille\\
edge processing & traitement local\\
packet & paquet de données\\
energy harvesting & récupération d'énergie\\
firmware update & mise à jour du micrologiciel\\
interoperability & interopérabilité\\
encryption & chiffrement
\end{tabularx}
\end{vocabulary}

\begin{applicationexercise}
A connected sensor is powered by a 2400 mAh battery. During each hour, it remains in deep sleep for
3590 s at a relative current of 2, then transmits for 10 s using LoRaWAN at a relative current of 18.
Assume that one relative-current unit corresponds to 0.10 mA.
\begin{enumerate}
\item Calculate the average current over one hour.
\item Estimate the ideal battery life in hours and in years.
\item Explain why the real service life would be shorter.
\item Suggest one change to the sensing or communication strategy that would extend autonomy.
\end{enumerate}
\end{applicationexercise}

\begin{questions}
\item Describe the information chain of an IoT system.
\item Why is edge processing useful for a battery-powered sensor?
\item Compare Wi-Fi, BLE and LoRaWAN for one application.
\item Why are software updates a safety and security requirement?
\item Give two methods for increasing sensor autonomy.
\end{questions}

\begin{engineeringchallenge}
Design a connected sensor for monitoring humidity in a school building. Select the sensing period,
communication method and power source. Define one alarm rule, one local fallback function and two
measures that protect user data.
\end{engineeringchallenge}

\begin{speaking}
Should every household appliance be connected? Present one useful service, one environmental cost and
one cybersecurity risk.
\end{speaking}
\end{multicols}

\subsection*{Sources}
\ECsource{NIST -- Internet of Things}{https://www.nist.gov/programs-projects/nist-cybersecurity-iot-program}
\ECsource{LoRa Alliance}{https://lora-alliance.org/}
